% Secondary figure: `paired-scatter` (design-axes.md axis 10; P4's frontier
% grammar). One point per condition: the strongest other arm on x, the
% primary arm on y, with the diagonal. Points above the diagonal (for a
% higher-is-better metric) favor the primary arm. Tier B from
% data/opening-pairs.dat written by make_paper_data.py --opening.
%
% source-data: runs/aggregate__*.json -> paper/data/opening-pairs.dat
% generator:   scripts/make_paper_data.py --opening
\begin{figure}[!t]
  \centering
  \acfitcol{%
  \begin{tikzpicture}
    \begin{axis}[acplot, width=0.62\linewidth, height=0.62\linewidth,
      xlabel={strongest other arm, \OPmetric}, ylabel={\OPprimarylabel, \OPmetric},
      xmin=\OPymin, xmax=\OPymax, ymin=\OPymin, ymax=\OPymax,
      legend style={at={(0.03,0.97)},anchor=north west,legend columns=1},
      \OPpointstyle,
    ]
      \addplot[ACMuted,dashed,forget plot,domain=\OPymin:\OPymax] {x};
      \addplot[only marks,mark=*,ACBlue,mark size=2pt] table[x=base,y=primary,meta=cond] {data/opening-pairs.dat};
      \legend{per-condition pairs}
    \end{axis}
  \end{tikzpicture}}
  \caption{\textbf{\slot{name what the pairing shows}.} Each point is one
    condition, plotting \OPprimarylabel\ against the strongest other arm in
    that condition; the dashed line is parity. \slot{one sentence stating
    how many points lie on the favorable side and which do not}.}
  \label{fig:distribution}
\end{figure}
